\section{Weekly Summary}

\begin{tcolorbox}[colback=yellow!10,colframe=black!50,title=AI-Generated Content]
\noindent The summary below was written by Claude (Anthropic) from that week's regression outputs and series descriptions. It has not been reviewed or fact-checked by a human, and should be read as an automated, informal recap -- not analysis.
\end{tcolorbox}

Welcome back to the weekly ritual of shoving two economic strangers into a room together and demanding they explain their relationship. As always, the machine paired these series at random, no thought or theory involved, so please adjust your expectations to "carnival fortune teller" levels before we begin. Spurious correlation isn't a risk here; it's the whole business model.

We opened Monday with Other Business Receivables Owned by Finance Companies versus a nonfarm diffusion index, and honestly, the two of them barely made eye contact. An R-squared of less than one percent and a p-value hovering around 0.10 means the raw regression already shrugged, and the detrended version shrugged in exactly the same way. No trend mirage here because there was nothing to mirage in the first place -- just two series politely ignoring each other. Tuesday was even more of a non-event, because pairing a monthly auto-production seasonal factor with the annual count of reissue patents originating in Romania produced a big pile of NaNs, which is the statistical equivalent of the computer refusing to dignify the question with an answer. Frankly, relatable.

Then things got spicy. Wednesday's matchup of the San Francisco Tech Pulse against the labor force of women with master's degrees actually got MORE interesting after detrending -- the R-squared climbed from about 10 percent to 14 percent, meaning this relationship survived having its shared drift surgically removed. That's the rare, respectable kind of result this project usually mocks, so I'll allow it a single dignified nod. Thursday's Employment Cost Index versus Food-at-Home CPI was the real drama queen: a jaw-dropping 95 percent R-squared in the raw regression, the kind of number that makes you want to write a manifesto. Except that's textbook shared-trend theater -- two things that both march upward over decades will always look like soulmates. Once detrended, R-squared collapsed to 43 percent. Still significant, sure, but "wages and grocery prices both track inflation" is roughly the least shocking sentence ever uttered, so let's not throw a parade.

Friday's Credit Intermediation PPI versus department store retail sales was another mutual snub, insignificant both raw and detrended, which is at least honest. And Saturday closed us out with Far West health care GDP versus portfolio investment income posting a gorgeous 97 percent raw R-squared that promptly deflated to 18 percent after detrending -- the classic "we only looked in love because we were both getting bigger over time" arc. The relationship technically stayed significant, but barely, and I wouldn't build a retirement plan on it. As ever, none of this means anything, no causal claims are being made, and if two random FRED series can fake a 97 percent correlation on trend alone, you should treat every headline correlation you ever see with the same suspicious side-eye. See you next week for more numerical speed dating.
